\documentclass[11pt]{article}

\usepackage[a4paper,margin=30mm]{geometry}
\usepackage{amsmath,amssymb,amsthm,mathtools}
\usepackage{enumitem}
\usepackage{microtype}
\usepackage{xcolor}
\usepackage[colorlinks=true,linkcolor=blue!55!black,citecolor=blue!55!black,
  urlcolor=blue!55!black]{hyperref}

\newtheorem{theorem}{Theorem}[section]
\newtheorem{proposition}[theorem]{Proposition}
\newtheorem{lemma}[theorem]{Lemma}
\newtheorem{corollary}[theorem]{Corollary}
\theoremstyle{definition}
\newtheorem{definition}[theorem]{Definition}
\theoremstyle{remark}
\newtheorem{remark}[theorem]{Remark}

\newcommand{\C}{\mathbb C}
\newcommand{\R}{\mathbb R}
\newcommand{\N}{\mathbb N}
\newcommand{\cS}{\mathcal S}
\newcommand{\XiR}{\Xi}
\newcommand{\dd}{\,\mathrm d}
\newcommand{\ord}{\operatorname{ord}}
\newcommand{\Rea}{\operatorname{Re}}
\newcommand{\Ima}{\operatorname{Im}}
\newcommand{\Diff}{\operatorname{Diff}}
\newcommand{\Mat}{\operatorname{Mat}}
\newcommand{\eqtag}[1]{\qquad\text{\upshape(#1)}}

\title{\textbf{The orientation-carrier boundary of the Riemann theta source}\\
\large Divisor width, Hilbert-metric collapse, and
hypertranscendental charge escape}
\author{Michael Haufschild}
\date{}

\begin{document}
\maketitle

\begin{abstract}
Put \(\Xi(z)=\xi(\frac12+iz)\), and retain the unsquared transverse
coordinate in the theta representation of \(\Xi\).  Its normal current is
\[
  J_a(b):=\frac{\partial}{\partial a}|\Xi(b-ia)|^2
  =2|\xi(\tfrac12+a+ib)|^2
     \Rea\frac{\xi'}{\xi}(\tfrac12+a+ib),
\]
with a zero-safe product form.  For every \(\alpha\geq0\), we prove
\(J_a(b)\geq0\) for all \(\alpha<a<1\) and \(b\in\R\) if and only if
every nontrivial zero \(\rho\) satisfies
\[
  \big|\Rea\rho-\tfrac12\big|\leq\alpha.
\]
The same threshold governs every complete Gaussian-windowed Hankel tower
and is independent of temperature.  Thus the current sign measures the
exact horizontal width of the Riemann divisor.

We classify natural carriers for this orientation.  At no positive
aperture can \(J_a\) factor as
\(|\Xi(b-ia)|^2Q(b)\) with \(Q\) entire.  Consequently the first
variation of the actual tilted theta autocorrelation has no finite Gram
representation by constant-coefficient differential filters.  The
global obstruction is stronger.  We prove a cyclic-null collapse law and
apply it to classify every jointly Schwartz-continuous, nonnegative,
momentum-symmetric form annihilating the theta source as a positive
Fourier measure on the real \(\Xi\)-divisor.  Every nonzero such form has
dense nonclosed ordinary-\(L^2\) radical, is nonclosable, and cannot
induce a positive-definite metric on the full theta-operator image.

The canonical unsquared Julia defect has the opposite polarity: every
off-critical quartet makes its positive first derivative carry
\(8\pi\) of Ward mass.  Functional-equation-preserving polynomial
surgeries insert arbitrary multiplicity \(N\) while leaving the complete
continuum real-aperture differential matroid unchanged.  That matroid
cannot recover \(N\); the source-marked noncompact boundary recovers
\(8\pi N\), even as its density escapes every bounded window.

This does not prove the Riemann hypothesis.  It closes the regular
source-core Hilbert carrier class for the theta-spectral program and
proves that canonical positive Julia orientation counts off-line defects
rather than excluding them.
\end{abstract}

\medskip
\noindent\textbf{Keywords.}
Riemann hypothesis; Riemann \(\xi\)-function; Jacobi theta function;
Hilbert--P\'olya; hypertranscendence; differential algebra; Ward charge;
Hankel moment problem.

\section{Introduction}

The Hilbert--P\'olya idea asks for a self-adjoint structure whose spectral
reality forces the nontrivial zeros of the Riemann zeta function onto the
critical line.  Hedenmalm recently obtained a concrete second-order
theta differential equation whose boundary eigenvalues are precisely the
zeros of
\[
  \Xi(z)=\xi\left(\frac12+iz\right),
  \qquad
  \xi(s)=\frac12s(s-1)\pi^{-s/2}\Gamma(s/2)\zeta(s),
\]
and isolated the missing step as the construction of an appropriate
sesquilinear structure for the operator pair
\cite{Hedenmalm2026}.  That construction raises a prior question.
What information must a source-side form retain before it can orient a
zero transversely?

The distinction matters because several natural operations preserve the
location or multiplicity of a divisor while erasing the direction from
which it was reached.  Squaring a normal coordinate is one such
operation.  Passing from a source to a scalar explicit formula is another.
An unsigned index can count an off-line pair without providing a
mechanism that makes the count vanish.

This paper keeps the normal coordinate unsquared.  The resulting
theta current
\[
  J_a(b)=\partial_a|\Xi(b-ia)|^2                     \eqtag{1.1}
\]
is the Fourier transform of the first variation of an autocorrelation of
the exponentially tilted Jacobi-theta source.  Its sign has a complete
geometric meaning: the first aperture beyond which the current is
nonnegative is exactly the horizontal half-width of the zero divisor.
This is stronger than another endpoint equivalence with the Riemann
hypothesis.  It determines the full quantitative phase diagram for every
\(\alpha\geq0\).

The positive result exposes a second problem.  A Hilbert-space argument
would naturally try to write \(J_a\) as a sum of squares.  We prove that
no finite local differential frame can do so.  In fact, the obstruction
is stronger: after division by \(|\Xi|^2\), the multiplier cannot extend
to any entire function.  Its divisor poles are necessary.  The proof
does not use simplicity of zeros.

We then close the corresponding regular global source-core problem for
Hedenmalm's theta pair.  Positivity and pair symmetry force every
jointly Schwartz-continuous form to be a Fourier measure on the real
\(\Xi\)-divisor.  The same theta source is cyclic in ordinary \(L^2\);
hence every nonzero form has a dense nonclosed radical and is
nonclosable.  No positive-definite form on the full algebraic image of
the theta operator can have such a pullback.  Thus passing from a finite
local frame to a closed or unbounded ordinary-Hilbert metric does not
repair the carrier.

The third result comes from hypertranscendental geometry.  For a real
aperture \(h\), the prime source
\[
  n\longmapsto \Lambda(n)(1-n^{-h})
\]
has a meromorphic Dirichlet carrier.  Distinct nonzero apertures give the
maximal possible differential matroid: zero aperture is a loop, repeated
apertures are parallel, and distinct nonzero apertures are differentially
independent over the rational coordinate field.  Nevertheless, an arbitrary number of off-critical
functional-equation quartets can be inserted without changing that
matroid.  The complete continuum aperture geometry is therefore
charge-blind.

The missing integer reappears at a source-marked boundary.  Each inserted
quartet produces a nonnegative pair of Lorentzians with mass \(8\pi\).
Multiplicity \(N\) gives mass \(8\pi N\).  Moving the quartet to
infinite height makes this density converge uniformly to zero on every
bounded window, while its total mass remains fixed.  The phenomenon is
not loss of charge but escape of charge.

This last positivity has the opposite logical direction from the theta
current.  The theta-current inequality would exclude off-line zeros;
the Julia--Ward inequality is already true for each off-line packet and
counts it.  Unsquaring preserves orientation in both constructions, but
does not choose its polarity.

\subsection{Relation to known instability and monotonicity results}

Horizontal monotonicity of \(|\xi|\) in zero-free half-planes is
classical and was used by Sondow and Dumitrescu to reformulate the
Riemann hypothesis \cite{SondowDumitrescu2010}; related monotonicity
criteria were developed by Matiyasevich, Saidak, and Zvengrowski
\cite{MatiyasevichSaidakZvengrowski2014}.  Our width law identifies the
entire phase transition of the concrete theta autocorrelation current and
its complete Gaussian Hankel hierarchy.  We do not claim novelty for
horizontal monotonicity in isolation.

The source-core metric theorem answers a different part of Hedenmalm's
proposal \cite{Hedenmalm2026}.  Its two harmonic-analysis inputs are
classical: Bochner--Schwartz represents positive-definite tempered
distributions by positive temperate measures
\cite{Schwartz1966}, and Wiener characterizes cyclic translation
vectors by almost-everywhere nonvanishing of their Fourier transform
\cite{Wiener1932}.  The new source-specific collision is that
\(\widehat k=\Xi\) activates both in opposite directions: theta nullity
forces the metric onto the real divisor, while the divisor's
Lebesgue-nullity makes the same source cyclic in ordinary \(L^2\).
Neither classical theorem alone gives the divisor classification,
dense radical, or nonclosability of Hedenmalm's pair.

Functional-equation-preserving perturbations are also known.
Gauthier and Xarles proved that failure of the corresponding Riemann
hypothesis is open and dense among broad spaces of meromorphic or entire
functions satisfying a fixed functional equation
\cite{GauthierXarles2009}.  Their proof already uses symmetric polynomial
interpolation to insert off-line zeros while approximating on compact
sets.  Accordingly, neither polynomial quartet insertion nor
compact-open instability is claimed here as new.  Their role is to
separate the new observables: continuum differential matroid, local
spectral carrier, and quantized marked boundary.

The differential-algebraic component belongs to the line beginning with
Ostrowski's work on Dirichlet series
\cite{Ostrowski1920,HansenStoneWolfson2022}; modern formulations include
Pong \cite{Pong2016}.  Difference analogues of logarithmic-derivative
estimates are much broader than the elementary zero-translation
argument used below \cite{HalburdKorhonen2006,ChiangFeng2009}.
The theorem package here is the coupling of these classical mechanisms
to the actual theta current and its source-marked boundary.

\subsection{What is decided}

The paper proves the following strict separation.
{\small
\[
\begin{array}{ccl}
\text{unsquared theta current}
  &\Longleftrightarrow& \text{horizontal divisor width},\\[1mm]
\text{finite local entire carrier}
  &\Longrightarrow& \text{impossible zero translation},\\[1mm]
\text{regular source-core Hilbert metric}
  &\Longrightarrow& \text{nonclosable divisor rigging},\\[1mm]
\text{aperture differential matroid}
  &\not\Longrightarrow& \text{quartet charge},\\[1mm]
\text{positive Julia--Ward boundary}
  &\Longrightarrow& \text{defect charge, not zero exclusion}.
\end{array}                                                    \eqtag{1.2}
\]
}
This is a classification inside explicitly stated observable classes,
not a no-go theorem for every Hilbert--P\'olya construction.  It also
does not prove that the Riemann divisor has width zero.  What it closes
is the information question: a proof based on the present theta source
cannot obtain zero exclusion from finite local Gram positivity,
an ordinary closed source-core Hilbert metric, compact-open control of
all jets, differential-algebraic maximality, or positivity of the
canonical Julia defect.  Those mechanisms are now classified rather
than left as unnamed candidates.

\section{Theta source, current, and aperture width}

\subsection{The normalized theta source}

Riemann's integration by parts gives the explicit real, even Schwartz
source
\[
 k(x)=\sum_{n=1}^{\infty}
 2\pi e^{x/2}
 \bigl(2\pi n^4e^{4x}-3n^2e^{2x}\bigr)
 e^{-\pi n^2e^{2x}},                                \eqtag{2.1}
\]
whose Jacobi reflection is \(k(-x)=k(x)\) and whose Fourier transform is
\[
  \Xi(b)=\int_{\R} k(x)e^{-ibx}\,\dd x.              \eqtag{2.2}
\]
For \(|a|<1\), define the
right-line source and tangent
\[
  k_a(x)=e^{-ax}k(x),\qquad \dot k_a(x)=-xk_a(x).
                                                               \eqtag{2.3}
\]
Both are Schwartz functions, and Fourier--Laplace continuation gives
\[
  \widehat{k_a}(b)=
  \xi\left(\frac12+a+ib\right)=\Xi(b-ia).            \eqtag{2.4}
\]

For \(f,g\in\mathcal S(\R)\), write
\[
  \operatorname{Corr}(f,g)(x)
  =\int_{\R} f(y+x)\overline{g(y)}\,\dd y.
\]
The first autocorrelation variation
\[
  A_a=\operatorname{Corr}(\dot k_a,k_a)
       +\operatorname{Corr}(k_a,\dot k_a)            \eqtag{2.5}
\]
has Fourier transform
\[
  \widehat{A_a}(b)=J_a(b)
  :=\frac{\partial}{\partial a}|\Xi(b-ia)|^2.        \eqtag{2.6}
\]
We call \(J_a\) the theta normal current.

\subsection{Zero-safe orientation identity}

\begin{proposition}[Normal orientation before scalarization]
\label{prop:current}
For \(|a|<1\) and \(b\in\R\),
\[
  J_a(b)
   =2\Rea\!\left[
       \xi'\!\left(\tfrac12+a+ib\right)
       \overline{\xi\!\left(\tfrac12+a+ib\right)}
     \right].                                       \eqtag{2.7}
\]
Equivalently, with \(s=\frac12+a+ib\),
\[
  J_a(b)
   =2|\xi(s)|^2\Rea\frac{\xi'}{\xi}(s)               \eqtag{2.8}
\]
whenever \(\xi(s)\neq0\).  The multiplicative identity obtained by
cross-multiplying (2.7) remains valid when \(\xi(s)=0\).
\end{proposition}

\begin{proof}
Differentiate
\(
|\Xi(b-ia)|^2
=\xi(s)\xi(\bar s)
\)
with respect to \(a\), use Schwarz symmetry
\(\xi(\bar s)=\overline{\xi(s)}\), and collect conjugate terms.
Division by \(|\xi(s)|^2\) gives (2.8) away from the divisor.  At a
zero, both sides of the cross-multiplied identity vanish.
\end{proof}

\begin{definition}[Horizontal divisor width]
\label{def:width}
For \(\alpha\geq0\), let
\[
  Z(\alpha):
  \qquad
  \xi(\rho)=0\Longrightarrow
  \left|\Rea\rho-\frac12\right|\leq\alpha.            \eqtag{2.9}
\]
The reflection symmetry \(\rho\mapsto1-\rho\) makes this equivalent to
the one-sided bound \(\Rea\rho\leq\frac12+\alpha\).
\end{definition}

\section{The exact aperture--width law}

\begin{theorem}[Normal-current phase law]
\label{thm:width}
For every \(\alpha\geq0\), the following are equivalent.
\begin{enumerate}[label=\textup{(\roman*)}]
\item \(Z(\alpha)\) holds.
\item For every \(a\) with \(\alpha<a<1\) and every \(b\in\R\),
      \[
        J_a(b)\geq0.                                 \eqtag{3.1}
      \]
\end{enumerate}
\end{theorem}

\begin{proof}
We first prove \(\textup{(i)}\Rightarrow\textup{(ii)}\).
Use the centered spectral coordinate
\[
  z=b-ia,\qquad \lambda=u+iv,
\]
where a zero \(\rho\) of \(\xi\) corresponds to a zero \(\lambda\)
of \(\Xi\) and \(|v|=|\Rea\rho-\frac12|\).  Pair each relative
resolvent atom with its conjugate.  Direct calculation gives
\[
\begin{aligned}
&\Ima\left[
 \left(\lambda^{-1}-(\lambda-z)^{-1}\right)
 +\left(\bar\lambda^{-1}-(\bar\lambda-z)^{-1}\right)
\right]\\
&\quad =
\frac{-2\,\Ima z\,
 \big((\Rea z-u)^2+(\Ima z)^2-v^2\big)}
 {\big((\Rea z-u)^2+(\Ima z-v)^2\big)
  \big((\Rea z-u)^2+(\Ima z+v)^2\big)}.
\end{aligned}
\eqtag{3.2}
\]
Here \(\Ima z=-a<-\alpha\) and \(|v|\leq\alpha\).
The denominator is positive away from a zero, while
\[
  (\Rea z-u)^2+a^2-v^2>0.
\]
Thus every conjugate pair in the canonical relative-resolvent trace
has nonnegative imaginary part.  The normally convergent paired sum
therefore has nonnegative imaginary part.  The logarithmic derivative
identity for the canonical product of \(\Xi\), followed by
Proposition~\ref{prop:current}, gives \(J_a(b)\geq0\).

For the converse, suppose (ii) and let
\(\rho=d+\frac12+ib\) be a zero with \(d>\alpha\).
Nontrivial zeros lie in \(0<\Rea s<1\), so \(0<d<\frac12\).
In the \(z\)-coordinate, \(z_0=b-id\) is a zero of \(\Xi\).
Zeros of a nonzero entire function are isolated.  Choose
\(c\in(\alpha,d)\), sufficiently close to \(d\), such that
\[
  \Xi(b-ix)\neq0\qquad(c\leq x<d).                   \eqtag{3.3}
\]
Set \(U(x)=|\Xi(b-ix)|^2\).  By (ii),
\(
U'(x)=J_x(b)\geq0
\)
on \([c,d]\), hence \(U(c)\leq U(d)\).  But (3.3) gives
\(U(c)>0\), while \(U(d)=0\), a contradiction.  Therefore no zero
has displacement \(d>\alpha\).  Reflection supplies the opposite
side, proving \(Z(\alpha)\).
\end{proof}

\begin{remark}
At \(\alpha=0\), Theorem~\ref{thm:width} specializes to a
monotonicity criterion equivalent to the Riemann hypothesis.  The
content used in this paper is the full law for every \(\alpha\), which
identifies what the current measures even when its endpoint is unknown.
\end{remark}

\subsection{The complete Gaussian Hankel phase transition}

Fix \(\tau>0\).  For each admissible aperture \(a\), define moments
\[
  m_n(a,\tau)
  =\int_{\R}b^nJ_a(b)e^{-\tau b^2}\,\dd b,\qquad n\geq0,            \eqtag{3.4}
\]
and the Hankel matrices
\[
  H_r(a,\tau)=\big(m_{j+k}(a,\tau)\big)_{0\leq j,k\leq r}.          \eqtag{3.5}
\]

\begin{corollary}[Temperature-independent Hankel threshold]
\label{cor:hankel}
Let \(\tau>0\) and \(\alpha\geq0\).  The following are equivalent:
\begin{enumerate}[label=\textup{(\roman*)}]
\item \(Z(\alpha)\);
\item for every \(\alpha<a<1\) and every \(r\geq0\),
      \(H_r(a,\tau)\) is positive semidefinite.
\end{enumerate}
Consequently, for any \(\tau,\sigma>0\), the two complete Hankel
tail conditions have exactly the same aperture threshold.
\end{corollary}

\begin{proof}
If \(J_a\geq0\), then for every polynomial \(p\),
\[
 \sum_{j,k}c_j\bar c_km_{j+k}(a,\tau)
 =\int_{\R}|p(b)|^2J_a(b)e^{-\tau b^2}\,\dd b\geq0,                \eqtag{3.6}
\]
so every Hankel matrix is positive semidefinite.

Conversely, positivity of the complete Hankel family says that the
Gaussian-weighted moment functional associated with the fixed continuous
density \(J_a\) is positive on polynomial squares.  Put
\(\dd\nu(b)=|J_a(b)|e^{-\tau b^2}\dd b\).  The standard vertical-strip
bounds for \(\xi\) and \(\xi'\) give a Gaussian tail for \(\nu\);
therefore polynomials are dense in \(L^2(\nu)\).  If \(J_a(b_0)<0\),
continuity gives a compact interval on which it is strictly negative.
Choose a nonzero continuous bump \(g\) supported in that interval and
polynomials \(p_n\to g\) in \(L^2(\nu)\).  Cauchy--Schwarz gives
\[
 \int_{\R}\bigl||p_n|^2-|g|^2\bigr|\,\dd\nu\longrightarrow0.
\]
The left side of (3.6) therefore converges to
\(\int |g|^2J_a(b)e^{-\tau b^2}\dd b<0\), a contradiction.  Hence
\(J_a\geq0\).
Theorem~\ref{thm:width} now proves the equivalence.  Since \(Z(\alpha)\)
does not involve \(\tau\), the threshold is temperature-independent.
\end{proof}

\section{Why regular Hilbert positivity cannot orient the current}

\subsection{Finite local carriers}

The phase law makes a sum-of-squares representation tempting.  We now
prove that the most direct local construction is impossible.

\begin{definition}[Entire diagonal carrier]
For \(a>0\), an entire diagonal carrier is an entire function
\(Q:\C\to\C\) satisfying, for every \(b\in\R\),
\[
\begin{aligned}
&\xi'(\tfrac12+a+ib)\xi(\tfrac12+a-ib)
 +\xi(\tfrac12+a+ib)\xi'(\tfrac12+a-ib)\\
&\qquad =
Q(b)\xi(\tfrac12+a+ib)\xi(\tfrac12+a-ib).            \eqtag{4.1}
\end{aligned}
\]
\end{definition}

\begin{theorem}[Entire-carrier obstruction]
\label{thm:entire}
No positive aperture \(a\) admits an entire diagonal carrier.
\end{theorem}

\begin{proof}
Assume that \(Q\) satisfies (4.1).  Both sides are entire functions of
the complex variable \(b\), and equality on \(\R\) implies equality on
\(\C\).  Put \(c=\frac12+a\) and substitute
\[
  b=-i(z-c).
\]
Using \(\xi(1-w)=\xi(w)\) and then dividing away from the divisor gives
\[
 \frac{\xi'}{\xi}(z)
 -\frac{\xi'}{\xi}(z-2a)
 =Q\big(-i(z-c)\big).                                \eqtag{4.2}
\]
The right side is entire.

Let \(\rho\) be any zero of \(\xi\), of arbitrary multiplicity \(m\).
The logarithmic derivative has residue \(m\) at \(\rho\).  Since the
right side of (4.2) has no pole, the pole at \(\rho\) must be cancelled
by the second term.  Therefore \(\rho-2a\) is also a zero, with the same
multiplicity.  Iteration gives
\[
  \rho-2na\in Z(\xi)\qquad(n=0,1,2,\ldots).           \eqtag{4.3}
\]
For large \(n\), the real part of \(\rho-2na\) is negative.  This
contradicts the fact that every nontrivial zero of the completed
\(\xi\)-function lies in \(0<\Rea s<1\).  Hence \(Q\) cannot exist.
\end{proof}

\begin{corollary}[Finite local autocorrelation-frame obstruction]
\label{cor:frame}
Let \(a\in(0,1)\).  There do not exist finitely many polynomials
\(p_1,\ldots,p_m\in\C[X]\) such that
\[
  A_a
  =\sum_{j=1}^m
    \operatorname{Corr}\big(p_j(\partial_x)k_a,
                            p_j(\partial_x)k_a\big),               \eqtag{4.4}
\]
where \(A_a\) is the actual theta autocorrelation tangent in (2.5).
\end{corollary}

\begin{proof}
Fourier transformation sends \(p_j(\partial_x)\) to the polynomial
symbol \(p_j(ib)\) and autocorrelation to modulus square.  Thus (4.4)
would give
\[
  J_a(b)=|\Xi(b-ia)|^2\sum_{j=1}^m|p_j(ib)|^2.       \eqtag{4.5}
\]
For real \(b\), the multiplier in (4.5) is the restriction of the
polynomial
\[
  Q(z)=\sum_{j=1}^m
  p_j(iz)\,\overline{p_j(i\bar z)}.
\]
Equation (4.5), written without division by \(\Xi\), is exactly
(4.1).  Theorem~\ref{thm:entire} gives the contradiction.
\end{proof}

\begin{remark}[What the obstruction says]
Infinite dimension alone does not evade Theorem~\ref{thm:entire} if the
resulting diagonal multiplier still extends holomorphically across the
divisor.  Any successful positive carrier in this theta model must retain
a non-entire transverse divisor singularity.  The theorem does not
exclude nonlocal or singular Hilbert structures; it proves that such a
departure is necessary.
\end{remark}

\subsection{The full translation-stable source core}

The preceding obstruction concerns diagonal multipliers.  We now allow
an arbitrary jointly continuous positive form on the Schwartz source
core.  This includes nonlocal and distributional kernels.

We first isolate the topological mechanism, which is independent of
zeta functions.

\begin{proposition}[Cyclic-null collapse principle]
\label{prop:cyclic-null}
Let \(\mathcal D\) be a complex vector space, let
\(j:\mathcal D\to H\) be an injective linear map into a Hilbert space,
and let \(Q\) be a Hermitian sesquilinear form on \(\mathcal D\).
Suppose that a family \((v_i)_{i\in I}\) lies in the left radical of
\(Q\) and that
\[
 \overline{\operatorname{span}\{j(v_i):i\in I\}}=H.               \eqtag{4.6}
\]
Then for every \(u\in\mathcal D\) there is a sequence
\((z_n)\subset\mathcal D\) such that
\[
\begin{gathered}
 j(z_n)\longrightarrow0,\qquad
 Q(z_m-z_n,z_m-z_n)=0,\\
 Q(z_n,z_n)=Q(u,u)\qquad(m,n\geq0).                \eqtag{4.7}
\end{gathered}
\]
Consequently, if \(Q\neq0\), its radical image is dense, proper, and
nonclosed in \(H\), and \(Q\) is not sequentially closable through
\(j\).
\end{proposition}

\begin{proof}
Choose \(r_n\) in the algebraic span of the \(v_i\) with
\(j(r_n)\to j(u)\), and put \(z_n=u-r_n\).  Hermitian symmetry makes
the left radical two-sided.  Expanding \(Q\) therefore gives (4.7).
The same cyclicity makes the radical image dense.  It is proper when
\(Q\neq0\), since injectivity of \(j\) would otherwise put every vector
in the radical.  A proper dense subspace is nonclosed.  Finally, (4.7)
is a form-Cauchy sequence converging to zero in \(H\), but its energy is
the fixed nonzero number \(Q(u,u)\) whenever such a \(u\) is chosen.
This is exactly the failure of sequential closability.
\end{proof}

Put
\[
 \Phi(x)=-\log k(x),\qquad
 D=-i\partial_x,\qquad
 L=-i(\partial_x+\Phi'(x))=kDk^{-1}.                 \eqtag{4.8}
\]
Then \(Lk=0\) and \(\ker L=\C k\) on the smooth source class.
Hedenmalm's pair-symmetry problem pulls back to
\[
 Q(Df,g)=Q(f,Dg),\qquad Q(k,k)=0.                    \eqtag{4.9}
\]
Here \(L\cS(\R)\) denotes only the algebraic image of the source;
no ordinary-\(L^2\) target domain is assumed.

\begin{theorem}[Source-core metric classification and collapse]
\label{thm:metric}
Let \(Q:\cS(\R)\times\cS(\R)\to\C\) be jointly continuous,
sesquilinear, and nonnegative, and suppose (4.9) holds.
Then:
\begin{enumerate}[label=\textup{(\roman*)}]
\item there is a unique positive Borel measure \(\mu\) of temperate
      growth such that
      \[
       Q(f,g)=\int_{\R}\widehat f(t)
              \overline{\widehat g(t)}\,\dd\mu(t);   \eqtag{4.10}
      \]
\item its support lies in the real Riemann divisor
      \[
       \operatorname{supp}\mu
       \subset\{t\in\R:\Xi(t)=0\};                   \eqtag{4.11}
      \]
      in particular \(\mu\) is purely atomic;
\item if \(Q\neq0\), its radical has dense, proper, nonclosed image
      in ordinary \(L^2(\R)\), and \(Q\) is not sequentially closable
      through \(\cS(\R)\hookrightarrow L^2(\R)\);
\item no nonzero form satisfying these hypotheses can have radical
      exactly \(\C k\).  Consequently it cannot be the pullback of a
      positive-definite form on the full algebraic image
      \(L\cS(\R)\).
\end{enumerate}
Conversely, every positive temperate measure supported on the set in
(4.11) defines a form satisfying the hypotheses.
\end{theorem}

\begin{proof}
Since \(D=-i\partial_x\), the first identity in (4.9) says that
\[
 Q(\partial_xf,g)+Q(f,\partial_xg)=0.
\]
Differentiating \(Q(T_yf,T_yg)\), where \(T_yf(x)=f(x+y)\), therefore
shows that \(Q\) is invariant under simultaneous translation.
The Bochner--Schwartz theorem gives the unique representation (4.10)
\cite{Schwartz1966}.

Nonnegativity supplies Cauchy--Schwarz for the seminorm of \(Q\), so
\(Q(k,k)=0\) puts \(k\) in both radicals.  Using
\(\widehat k=\Xi\) in (4.10) gives
\[
 0=Q(k,k)=\int_{\R}|\Xi(t)|^2\,\dd\mu(t).            \eqtag{4.12}
\]
The integrand is continuous and nonnegative, proving (4.11).
The real zeros of a nonzero entire function are discrete, hence every
such measure is atomic.  The converse follows directly from (4.10).

Translation invariance puts every \(T_yk\) in the radical.  Because
\(\Xi\) vanishes only on a Lebesgue-null discrete set, it is nonzero
almost everywhere on \(\R\).  Wiener's \(L^2\) cyclicity theorem
\cite{Wiener1932} gives
\[
 \overline{\operatorname{span}\{T_yk:y\in\R\}}^{\,L^2}
 =L^2(\R).                                          \eqtag{4.13}
\]
Proposition~\ref{prop:cyclic-null}, applied to the translation orbit,
now gives the dense proper nonclosed radical and nonclosability.
Finally,
positive definiteness on \(L\cS(\R)\) would make the pullback radical
equal to \(\ker L=\C k\).  Its \(L^2\)-image is one-dimensional and
closed, contradicting the dense-radical conclusion.
\end{proof}

\begin{remark}[Exact scope]
Theorem~\ref{thm:metric} classifies the natural translation-stable
Schwartz source extension.  A form defined only on a finite
zero-eigenfunction span, with no such extension, lies outside the
theorem.  The result is therefore a source-core no-go for Hedenmalm's
theta pair, not a no-go for de Branges spaces or Hilbert--P\'olya in
general.
\end{remark}

\section{Continuum differential geometry and quartet blindness}

\subsection{The real-aperture prime family}

For \(h\in\R\), define the arithmetic source
\[
  a_h(n)=\Lambda(n)(1-n^{-h})                         \eqtag{5.1}
\]
and, initially in a right half-plane, its Dirichlet carrier
\[
\begin{aligned}
  D_h(s)
  &=\sum_{n\geq1}\frac{\Lambda(n)(1-n^{-h})}{n^s}\\
  &=\frac{\zeta'}{\zeta}(s+h)-\frac{\zeta'}{\zeta}(s).              \eqtag{5.2}
\end{aligned}
\]
If
\[
 A(s)=\tfrac12s(s-1)\pi^{-s/2}\Gamma(s/2),
\]
then on the common Euler half-plane (5.2) is exactly
\[
 \left[\frac{\xi'}{\xi}(s+h)-\frac{\xi'}{\xi}(s)\right]
 -
 \left[\frac{A'}{A}(s+h)-\frac{A'}{A}(s)\right].     \eqtag{5.3}
\]
Thus \(D_h\) is the prime-corrected completed logarithmic shift defect;
the archimedean term has been subtracted explicitly, not absorbed into
the differential base.

For a family \(f_i\) in a differential extension of \(\C(s)\), let
\(\Mat(f_i)\) denote the represented differential matroid:
\(I\) is independent when \(\{f_i:i\in I\}\) is differentially
algebraically independent over \(\C(s)\).

\begin{lemma}[Prime-tail differential criterion]
\label{lem:ostrowski}
Let \(f_1,\ldots,f_r\) be Dirichlet series with finite abscissa of
convergence and coefficients supported on positive prime powers.  If,
for every \(m\), the prime-layer coefficient profiles of
\[
 f_j,f_j',\ldots,f_j^{(m)}\qquad(1\leq j\leq r)
\]
are linearly independent modulo functions supported on finitely many
primes, then no nonzero differential polynomial over \(\C(s)\)
vanishes on \(f_1,\ldots,f_r\).
\end{lemma}

\begin{proof}
This is the prime-support form of Ostrowski's coefficient criterion
\cite{Ostrowski1920,Pong2016}.  Expand a hypothetical
differential-polynomial relation as a Dirichlet series and choose its
lowest nonvanishing coefficient.  Differentiation multiplies the
\(n\)-th coefficient by powers of \(-\log n\), while multiplication
becomes Dirichlet convolution.  Restriction to new prime rays and
polarization of the highest nonzero homogeneous part produce a
finite-prime-tail linear relation among the displayed jet profiles.
The hypothesis kills that relation.  Descending through the homogeneous
parts makes the original differential polynomial zero.
\end{proof}

\begin{proposition}[Pointed-collision aperture matroid]
\label{prop:matroid}
For every finite real-aperture family \(h_i\),
\[
  \Mat(D_{h_i})
  =\text{the pointed-collision matroid of }(h_i):      \eqtag{5.4}
\]
\(h_i=0\) is a loop, equal nonzero apertures are parallel, and any
collection of distinct nonzero apertures is differentially algebraically
independent over \(\C(s)\).
\end{proposition}

\begin{proof}
The first two statements follow directly from (5.1).  For distinct
nonzero \(h_1,\ldots,h_r\), take all logarithmic jets.  On a prime \(p\),
their coefficient profiles are finite linear combinations of
\[
  (\log p)^{k+1},
  \qquad
  e^{-h_j\log p}(\log p)^{k+1}
  \quad(1\leq j\leq r,\ k\geq0).                     \eqtag{5.5}
\]
The real exponents \(0,h_1,\ldots,h_r\) are distinct.  The
exponential-polynomial Vandermonde lemma implies that no nontrivial
finite combination of the profiles in (5.5) vanishes on a tail of the
primes.  Lemma~\ref{lem:ostrowski} gives differential algebraic
independence.
\end{proof}

\subsection{Symmetric finite-divisor surgery}

Fix \(d\in(0,\frac12)\), \(B>0\), and put \(z=s-\frac12\).  Define
\[
  G_{d,B}(s)
  =\left(1+\frac{(z-d)^2}{B^2}\right)
   \left(1+\frac{(z+d)^2}{B^2}\right).               \eqtag{5.6}
\]
It has zeros
\[
  \frac12\pm d\pm iB,                                \eqtag{5.7}
\]
satisfies \(G_{d,B}(1-s)=G_{d,B}(s)\), has real
coefficients, and is positive on the real segment \(0<s<1\).
For \(N\in\N\), set
\[
  F_{N,B}(s)=G_{d,B}(s)^N\xi(s).                     \eqtag{5.8}
\]
For the escape statement below, \(B_n=n+1\).

\begin{proposition}[All-jet escape with arbitrary charge]
\label{prop:escape}
For each \(N\geq1\):
\begin{enumerate}[label=\textup{(\roman*)}]
\item \(F_{N,B}\) is entire, Schwarz-real, satisfies
      \(F_{N,B}(1-s)=F_{N,B}(s)\), is hypertranscendental over
      \(\C(s)\), and its polynomial gauge contributes the
      off-critical quartet (5.7) with multiplicity \(N\);
\item as \(n\to\infty\), every holomorphic derivative of \(F_{N,B_n}\)
      converges to the corresponding derivative of \(\xi\), locally
      uniformly on \(\C\).
\end{enumerate}
\end{proposition}

\begin{proof}
All symmetries and zeros follow from (5.6).  Multiplication by a nonzero
polynomial preserves hypertranscendence.  On every compact set,
\[
  G_{d,B}(s)
  =1+\frac{2(z^2+d^2)}{B^2}
    +\frac{(z^2-d^2)^2}{B^4}\longrightarrow1.        \eqtag{5.9}
\]
The convergence is locally uniform.  Holomorphic convergence and
Cauchy's integral formula give convergence of every derivative; fixed
powers and multiplication by \(\xi\) preserve it.
\end{proof}

Let
\[
  q_{d,B}(s)=\frac{G'_{d,B}(s)}{G_{d,B}(s)}.          \eqtag{5.10}
\]
For the completed shift defect, surgery gives the exact affine change
\[
 D_h^{(N,B)}(s)
 =D_h(s)+N\big(q_{d,B}(s+h)-q_{d,B}(s)\big),          \eqtag{5.11}
\]
after the same fixed archimedean correction on both sides.  The second
term is rational.

\begin{theorem}[Continuum hypertranscendental charge blindness]
\label{thm:blind}
For every aperture family \(h_i\), every \(B>0\), and every
\(N\in\N\),
\[
  \Mat(D_{h_i}^{(N,B)})=\Mat(D_{h_i}).                \eqtag{5.12}
\]
Hence there is no function from the complete real-aperture differential
matroid to \(\N\) that returns the inserted quartet multiplicity for all
\(N\).
\end{theorem}

\begin{proof}
Each correction in (5.11) belongs to the differential base
\(\C(s)\).  Translating a represented family by base-field elements
does not change any differential algebraic dependence relation:
\[
 \C(s)\langle D_{h_i}^{(N,B)}:i\in I\rangle
 =\C(s)\langle D_{h_i}:i\in I\rangle.                \eqtag{5.13}
\]
This proves (5.12).
If a matroid classifier returned \(N\), evaluate it at \(N=1\) and
\(N=2\).  The inputs are equal by (5.12), while the required outputs
are different, a contradiction.
\end{proof}

\begin{remark}
Theorem~\ref{thm:blind} is not the assertion that
hypertranscendence is weak.  Proposition~\ref{prop:matroid} says that
the aperture family has maximal differential rank.  The stronger fact
is that the whole rank geometry is transverse to the finite divisor
charge.
\end{remark}

\section{The marked noncompact boundary recovers the charge}

Set \(c=d+\frac12>0\).  The strict-right Schwarz pair in (5.7) has
the unsquared Julia phase defect
\[
 \Delta_{d,B,t}(\varepsilon)
 =1-\prod_{\eta\in\{\pm1\}}
 \frac{(c-\varepsilon)^2+(t+\eta B)^2}
      {(c+\varepsilon)^2+(t+\eta B)^2}.              \eqtag{6.1}
\]
For \(\varepsilon\geq0\), every factor in the product lies in
\([0,1]\), so \(\Delta_{d,B,t}(\varepsilon)\geq0\).
Direct differentiation before any divisor squaring gives
\[
 \left.\partial_\varepsilon
   \Delta_{d,B,t}(\varepsilon)\right|_{\varepsilon=0}
 =2\left[
 \frac{2c}{c^2+(t+B)^2}
 +\frac{2c}{c^2+(t-B)^2}
 \right].                                           \eqtag{6.2}
\]
Thus this canonical positive Julia orientation has a
\emph{defect-counting} polarity: an off-critical pair creates a
strictly positive first derivative.  It is not the inequality
\(J_a\geq0\) of Theorem~\ref{thm:width}, whose validity would exclude
zeros beyond the aperture.

The same pair has
the reflected Cauchy response
\[
 \mathcal R_{d,B}(w)
 =2\left[
   \frac1{w+d+iB}+\frac1{w+d-iB}
 \right].                                           \eqtag{6.3}
\]
In the normalization fixed by one logarithmic source pole, its
Euler-edge Ward density is
\[
 W_{d,B}(t)
 =2\left[
  \frac{2c}{c^2+(t+B)^2}
 +\frac{2c}{c^2+(t-B)^2}
 \right]
 =2\Rea\mathcal R_{d,B}\left(\frac12+it\right).      \eqtag{6.4}
\]
Equation (6.2) says exactly that
\(W_{d,B}(t)=\partial_\varepsilon
\Delta_{d,B,t}(0)\).
For multiplicity \(N\), put \(W_{d,B}^{(N)}=NW_{d,B}\).

\begin{theorem}[Quantized boundary recovery and escape]
\label{thm:ward}
Let \(d\in(0,\frac12)\), \(B>0\), and \(N\in\N\).
\begin{enumerate}[label=\textup{(\roman*)}]
\item \(W_{d,B}^{(N)}(t)\geq0\) for all \(t\), and
      \[
        \int_{\R}W_{d,B}^{(N)}(t)\,\dd t=8\pi N.      \eqtag{6.5}
      \]
      Thus total marked boundary mass is injective in \(N\).
\item For every \(R\geq0\) and \(B>R\),
      \[
        \sup_{|t|\leq R}W_{d,B}^{(N)}(t)
        \leq \frac{8N}{(B-R)^2}.                     \eqtag{6.6}
      \]
      Hence for fixed \(N\), the density tends uniformly to zero on
      every bounded height window as \(B\to\infty\), although its mass
      remains \(8\pi N\).
\item The marked logarithmic source \(Nq_{d,B}\) has a simple pole at
      every inserted point.  No nonzero constant-coefficient
      differential operator annihilates its principal-part class.
\end{enumerate}
\end{theorem}

\begin{proof}
Each term in (6.4) is a positive Lorentzian.  Translation invariance and
\[
  \int_{\R}\frac{2c}{c^2+t^2}\,\dd t=2\pi            \eqtag{6.7}
\]
give (6.5).  Since \(8\pi>0\), equality of the masses for two natural
numbers forces equality of the multiplicities.

If \(|t|\leq R<B\), then
\(
|t\pm B|\geq B-R
\)
and \(c<1\).  Dropping the positive \(c^2\) in each denominator gives
\[
 W_{d,B}(t)
 \leq2\left[\frac2{(B-R)^2}+\frac2{(B-R)^2}\right]
 =\frac8{(B-R)^2},
\]
which proves (6.6).

Finally, a logarithmic derivative of a nonzero analytic function has a
simple pole at each zero, independent of the zero's multiplicity as a
statement about pole order.  More precisely, if a local factor is
\((s-\rho)^Nm(s)\), \(m(\rho)\neq0\), then
\[
  Nq_{d,B}(s)=\frac{N}{s-\rho}+\frac{m'(s)}{m(s)}.    \eqtag{6.8}
\]
If \(P(\partial_s)\) is nonzero of degree \(r\), applying it to
\((s-\rho)^{-1}\) creates a pole of order \(r+1\) whose leading
coefficient is nonzero.  Lower derivatives cannot cancel it.
\end{proof}

\subsection{The singular common mode}

There is a useful exact form of the last statement.  Let
\(\mathcal M_\rho/\mathcal O_\rho\) be meromorphic germs at \(\rho\)
modulo holomorphic germs, and let \(I\) be a finite nonempty set of
source channels carrying the same marked logarithmic mode.  Define
\[
\begin{aligned}
\Phi_I:\C[\partial]^{I}
   &\longrightarrow \mathcal M_\rho/\mathcal O_\rho,\\
 (P_i)_{i\in I}
   &\longmapsto
\left[\sum_{i\in I}P_i(\partial)(Nq_{d,B})\right].
\end{aligned}                                                        \eqtag{6.9}
\]
Then
\[
 \ker\Phi_I
 =\left\{(P_i):\sum_{i\in I}P_i=0\right\},\qquad
 \operatorname{im}\Phi_I
 =\C[\partial]\,[Nq_{d,B}].                           \eqtag{6.10}
\]
Thus the principal-part source projection has one free marked mode, with
the \(A_{|I|-1}\) augmentation kernel.  This is the algebraic location
of the charge that the aperture matroid forgets.

\section{The orientation-carrier theorem}

We can now collect the results into one visibility law.

\begin{theorem}[Orientation-carrier fracture]
\label{thm:main}
The Riemann theta source has the following five exact layers.
\begin{enumerate}[label=\textup{\Alph*.}]
\item \textbf{Geometric visibility.}
      The sign phase of the unsquared normal current is exactly the
      horizontal width of the Riemann divisor
      (Theorem~\ref{thm:width}), equivalently the phase threshold of
      every complete Gaussian Hankel tower
      (Corollary~\ref{cor:hankel}).
\item \textbf{Local-carrier obstruction.}
      At every positive aperture the current has no entire diagonal
      carrier and the actual theta autocorrelation tangent has no
      finite constant-coefficient differential Gram frame
      (Theorem~\ref{thm:entire} and
      Corollary~\ref{cor:frame}).
\item \textbf{Regular global collapse.}
      Every nonnegative momentum-symmetric form on the
      translation-stable Schwartz source core is a singular measure on
      the real \(\Xi\)-divisor.  Every nonzero such form is nonclosable
      in ordinary \(L^2\), has dense nonclosed radical, and cannot
      induce a positive-definite form on the full theta-operator image
      (Theorem~\ref{thm:metric}).
\item \textbf{Polarity fracture.}
      The canonical Julia phase defect is nonnegative, but an
      off-critical quartet makes its unsquared first derivative
      strictly positive.  Its total derivative mass is \(8\pi N\);
      therefore this positivity counts defects rather than excluding
      them (equations (6.1)--(6.5)).
\item \textbf{Hypertranscendental blindness and boundary escape.}
      The complete continuum aperture family has the pointed-collision
      differential matroid, but this matroid is unchanged by an
      off-critical quartet of arbitrary multiplicity.  The surgery
      converges through every holomorphic jet on compact sets, while
      the source-marked noncompact boundary alone recovers the
      multiplicity as \(8\pi N\); its density nevertheless disappears
      on every bounded window as the quartet escapes
      (Proposition~\ref{prop:matroid} and
      Theorem~\ref{thm:blind},
      Proposition~\ref{prop:escape}, and
      Theorem~\ref{thm:ward}).
\end{enumerate}
\end{theorem}

\begin{proof}
Parts A--E are precisely the cited results.  Their conjunction is
nonvacuous: Part A proves that the current contains the desired
orientation; Part B shows that a finite local sum-of-squares carrier
cannot explain it; Part C closes the full regular ordinary-Hilbert
source-core extension; Part D proves that the canonical positive
unsquared alternative has the counting polarity; Part E proves that
even maximal differential rank and all compact-open jets miss the
charge that the marked noncompact boundary retains.
\end{proof}

\subsection{Consequence for theta-spectral RH programs}

Hedenmalm's theta equation provides eigenfunctions for all zeros, including
hypothetical nonreal spectral parameters, while a suitable positive
sesquilinear structure would force those parameters to be real
\cite{Hedenmalm2026}.  Theorem~\ref{thm:main} determines a necessary
carrier property for such a structure in the present source model:

\[
\boxed{
\begin{minipage}{0.86\linewidth}
Within the translation-stable Schwartz source model, the desired theta
orientation has neither a finite local entire Gram carrier nor a
nonzero closable ordinary-\(L^2\) positive-definite carrier.  The
canonical positive Julia carrier does retain an unsquared orientation,
but with the opposite polarity: it assigns positive quantized mass to
each off-line defect.
\end{minipage}}                                                     \eqtag{7.1}
\]

This goes beyond loss of sign under squaring.  The preserved-sign
current exists and exactly measures width; its
regular Hilbert carriers are excluded.  The positive Julia derivative
also exists; its sign is already satisfied by the unwanted divisor.
Compact-open control of every holomorphic jet and the complete aperture
differential matroid remain unable to distinguish multiplicity, while
the only observable in the theorem package that does distinguish it is
singular and noncompact.

This conclusion is also narrower than a no-go theorem for
Hilbert--P\'olya.  A genuinely global operator, a de Branges space, or a
trace formula may contain precisely the singular and noncompact
information isolated here.  The theorem says how such a construction can
be tested: if it factors through an observable proven charge-blind above,
it cannot force the critical line.

\section{Closed scope and conclusion}

The paper answers five questions without replacing them by conjectures.

\begin{enumerate}[label=\textup{(\arabic*)}]
\item \emph{Does keeping the Julia/theta coordinate unsquared preserve
      transverse orientation?}  Yes.  Its normal current is (2.6), and
      its sign threshold is exactly the full horizontal divisor width.
\item \emph{Can a finite local positive spectral frame supply that sign?}
      No.  Any entire diagonal carrier forces an impossible translation
      orbit of every zero.
\item \emph{Can a regular global ordinary-Hilbert source metric supply
      it?}  No.  Every nonzero positive pair-symmetric source-core form
      is a singular divisor measure with dense nonclosed radical and is
      nonclosable.
\item \emph{Does the canonical positive Julia derivative have the
      exclusion polarity?}  No.  It is strictly positive when an
      off-critical quartet is present and counts it in quanta of
      \(8\pi\).
\item \emph{Can complete hypertranscendental aperture geometry detect
      off-line multiplicity?}  No.  All multiplicities have the same
      continuum differential matroid; the missing multiplicity remains
      observable only in the source-marked noncompact boundary, whose
      exact mass is \(8\pi N\) and whose principal-part mode has kernel
      (6.10).
\end{enumerate}

No assertion in these answers says that the actual marked mass of the
Riemann divisor vanishes.  Such an assertion would prove the Riemann
hypothesis and is not obtained here.  The advance is a carrier theorem:
the actual theta source already contains an exact quantitative
orientation observable, and the analysis completely classifies the
local, regular source-Hilbert, compact-open, and differential-algebraic
carriers considered here.  It also resolves the apparent Julia remedy:
unsquaring restores a sign, but the canonical positive sign is a
quantized defect counter.  No unproved vanishing or tightness statement
is included among the conclusions.

\begin{thebibliography}{99}
\small

\bibitem{ChiangFeng2009}
Y.-M. Chiang and S.-J. Feng,
\newblock On the growth of logarithmic differences, difference quotients
and logarithmic derivatives of meromorphic functions,
\newblock \emph{Trans. Amer. Math. Soc.} \textbf{361} (2009),
3767--3791.

\bibitem{GauthierXarles2009}
P. M. Gauthier and X. Xarles,
\newblock Perturbations of \(L\)-functions with or without non-trivial
zeros off the critical line,
\newblock in \emph{New Directions in Value-Distribution Theory of Zeta
and \(L\)-Functions}, 2009; arXiv:0911.5135.

\bibitem{HalburdKorhonen2006}
R. G. Halburd and R. J. Korhonen,
\newblock Difference analogue of the lemma on the logarithmic derivative
with applications to difference equations,
\newblock \emph{J. Math. Anal. Appl.} \textbf{314} (2006), 477--487.

\bibitem{HansenStoneWolfson2022}
E. C. Hansen, Y. Stone, and J. Wolfson,
\newblock Alexander Ostrowski's ``On Dirichlet series and algebraic
differential equations'',
\newblock arXiv:2211.02088, 2022.

\bibitem{Hedenmalm2026}
H. Hedenmalm,
\newblock Spectral interpretation of Riemann zeta zeros,
\newblock arXiv:2606.17494, 2026.

\bibitem{MatiyasevichSaidakZvengrowski2014}
Y. Matiyasevich, F. Saidak, and P. Zvengrowski,
\newblock Horizontal monotonicity of the modulus of the Riemann zeta
function and related functions,
\newblock \emph{Acta Arith.} \textbf{166} (2014), 189--200.

\bibitem{Ostrowski1920}
A. Ostrowski,
\newblock \"Uber Dirichletsche Reihen und algebraische
Differentialgleichungen,
\newblock \emph{Math. Z.} \textbf{8} (1920), 241--298.

\bibitem{Pong2016}
W. Y. Pong,
\newblock Applications of differential algebra to algebraic independence
of arithmetic functions,
\newblock \emph{Acta Arith.} \textbf{172} (2016), 149--173.

\bibitem{SondowDumitrescu2010}
J. Sondow and C. Dumitrescu,
\newblock A monotonicity property of Riemann's xi function and a
reformulation of the Riemann hypothesis,
\newblock arXiv:1005.1104, 2010.

\bibitem{Schwartz1966}
L. Schwartz,
\newblock \emph{Th\'eorie des distributions},
\newblock Hermann, Paris, 1966.

\bibitem{Wiener1932}
N. Wiener,
\newblock Tauberian theorems,
\newblock \emph{Ann. of Math.} \textbf{33} (1932), 1--100.

\end{thebibliography}

\end{document}
